\chapter{研究背景与挑战}

\section{选题背景与意义}

根据世界卫生组织的统计，全世界有xxx人罹患残疾性听力损失。手语作为听障人士的主要交流语言，帮助听障人之间、听障人与健听人之间沟通，发挥着弥合鸿沟的作用。手语通过手部和肢体动作、面部表情传递信息，且和其他自然语言的语序有所区别，这给理解、学习手语带来了困难。

为了辅助听障人士更方便地交流，基于计算机视觉和深度学习的手语理解技术应运而生。通过分析手部、肢体动作和面部表情，手语理解算法能从视频中识别出对应的内容，从而辅助手语学习、理解等场景。

\section{研究挑战}

